%
%============= CHAPTER BREAK =================%
%
\chapter{Lucian of Samosata (145--185 CE)}%
\justifying%
\vspace*{3mm}%
%
\noindent{}A prolific satirist and essayist of the mid- to late-2nd century, %
Lucian is well-known to scholars of early Christianity, %
most especially for his roasting of the sometime Christian \textit{Peregrinus}, %
whose self-immolation Lucian witnessed during the Olympic games of 169 CE. \\[2mm] %

\noindent{}\textbf{προσδαπανάω} / \textit{prosdapanaō} / “to spend extra” %
has 90 TLG hits. %
The first known non-Christian use is most likely in Lucianʼs \textit{Saturnalia}, %
otherwise known as \textit{Cronosolon} / Κρονοσόλων.\endnote{%
There is an occurrence in one notable manuscript tradition (cod. Marcianus 408) %
of the \textit{History of Alexander the Great} (Reichmann 1963), %
but the term is absent from the older, overlapping manuscript traditions that %
may trace back to a mid-1st century CE textual form.} % 
As Popescu has described, the three distinct sections of the work treat of similar subject matter, %
while changing perspectives and roles.\endnote{Popescu 2016:219.} %
This “mini-corpus” contains a succession of dialogues, laws, and epistles; %
it thus enacts an eclectic, metatextual performance and commentary %
on Lucianʼs many previous, more substantive essays.\endnote{%
Popescu 2016:219--221.} %
In two of the dialogues, Lucian wears the guise of a priest of Cronus, %
but he also apparently depicts himself, by way of the fettered Cronus, %
as “an old, gout-ridden word-charmer”.\endnote{Popescu 2016:222.} %
This framing aligns at a general level with the workʼs date %
deduced by Fowler \& Fowler, around 165 CE.\endnote{%
Fowler \& Fowler 1905, 1:xvi.} 
The Fowlers additionally imply the coherence of the materials when %
describing the “Saturnalia” as “dialogue and letters on the relation of rich and poor”.\endnote{\textit{Idem}.} %

The new/rare lemma appears at the very end of the \textit{Saturnalia}, %
in a letter the rich collectively write back to Cronus, %
right after they have accused their impoverished guests of several gross improprieties. \clearpage

 \begin{quote}Ταῦτ’ ἐστὶ καὶ τὰ τοιαῦτα, 
 ὑφ’ ὧν ἡμεῖς ἐβουλευσάμεθα πρὸς τὸ λοιπὸν ἀσφαλείας τῆς ἡμετέρας %
 ἕνεκεν μηκέτι ἐπιβατὸν αὐτοῖς ποιεῖν τὴν οἰκίαν. % 
 εἰ δὲ ἐπὶ σοῦ συνθοῖντο μετρίων δεήσεσθαι, %
 ὥσπερ νῦν φασι, μηδὲν δὲ ὑβριστικὸν ἐν τοῖς συμποσίοις ἐργάσεσθαι, %
 κοινωνούντων ἡμῖν καὶ συνδειπνούντων τύχῃ τῇ ἀγαθῇ. %
 καὶ τῶν ἱματίων, ὡς σὺ κελεύεις, %
 πέμψομεν καὶ τοῦ χρυσίου ὁπόσον οἷόν τε καὶ \textbf{προσδαπανήσομεν}, %
 καὶ ὅλως οὐδὲν ἐλλείψομεν. %
 καὶ αὐτοὶ δὲ ἀφέμενοι τοῦ κατὰ τέχνην ὁμιλεῖν %
 ἡμῖν φίλοι ἀντὶ κολάκων καὶ παρασίτων ἔστωσαν. (\textit{Saturn}. 39)\endnote{%
 Page et al, 1959:138.} \\[1mm] %
 These things, and such things as these, %
 based on which we took counsel, %
 for the sake of what remains of our own security, %
 no longer to make the house accessible to them. %
 But if they should covenant with you to restrain themselves by limits, %
 just as they now say, %
 and participate in no insolence at the symposiums, %
 let them share with us and let them feast together with good fortune. %
 And as you command, %
 we will send them cloaks and gold, %
 as much as such, and \textbf{we will spend extra}, %
 and nothing at all will we lack. %
 And so now let them, abandoning their crafty consorting %
 with us, be friends instead of flatterers and parasites. \end{quote} % 
 \vspace*{2mm}

This new lemma also appears in Luke 10:35, %
in the famous fable of the Good Samaritan, %
a story attested as absent for Marcionʼs \textit{Evangelion}. %
After an anonymous traveler is beaten, left half-dead, %
and passed by a priest and a Levite, %
a Samaritan has compassion on him, %
tends to his wounds, puts him on his pack animal, %
and brings him to an inn. %
He then instructs the innkeeper: %

\vspace*{2mm} 
 \begin{quote}Ἐπιμελήθητι αὐτοῦ, %
 καὶ ὅ τι ἂν \textbf{προσδαπανήσῃς} %
 ἐγὼ ἐν τῷ ἐπανέρχ\-εσθαί με ἀποδώσω σοι. \\[1mm] %
 Take care of him, %
 and whatever \textbf{you should spend extra}, %
 I myself, on my journeying back, will give back to you. % 
 \end{quote} % 
\vspace*{2mm}

The connection is both deeper and broader than the signature lemma. %
In both cases, the treatment of the poor by the rich is in focus. %
Both texts use playful storytelling to illustrate the dilemma, %
both are stories within stories, %
and both follow after an ethical dialogue/debate. %
Both also accentuate the theme of hospitality to the poor. %
The Lukan story is framed by a question from a law-expert, %
“who is my neighbor?” %
Lucianʼs dialogues, laws, and letters in the \textit{Saturnalia} all focus on the obligation %
of the wealthy to give to the poor. %
Indeed, Lucianʼs law-giver (Cronosolon) elaborates %
the specific obligation of the wealthy to give “a tenth” / τὸ δέκατον / \textit{to dekaton} %
of their annual income to their poor neighbors (\textit{Saturn}. 14).\endnote{%
Page et al 1959:108. Poor writers got double during Saturnalia, O Patrons!} %

The precise, shared lemma and this larger array of stylistic and thematic %
similarities cannot be mere coincidence. %
The editor of canonical Luke was almost certainly inspired by the learned composition %
of Lucian on the Saturnalia, %
pushing the date of the canonical form of Luke after 165 CE.\endnote{%
On a passing note, for scholars such as Klinghardt who have researched %
meals and feasts in canonical Luke, %
there are many other parallels with Lucianʼs \textit{Saturnalia} that are worth exploring!} %